\section*{Chapitre \ref{ch:g12:trigo} --- Fonctions trigonométriques}

\begin{solution}{exo:g12:trigo:1}
$\cos\frac{2\pi}{3} = \cos\left(\pi - \frac\pi3\right) = -\cos\frac\pi3
= -\frac12$.

$\sin\frac{5\pi}{6} = \sin\left(\pi - \frac\pi6\right) = \sin\frac\pi6
= \frac12$.

$\cos\frac{7\pi}{4} = \cos\left(2\pi - \frac{\pi}{4}\right)
= \cos\left(-\frac\pi4\right) = \frac{\sqrt2}{2}$.

$\sin\left(-\frac\pi3\right) = -\sin\frac\pi3 = -\frac{\sqrt3}{2}$.
\end{solution}

\begin{solution}{exo:g12:trigo:2}
$f'(x) = 2\sin x\cos x = \sin 2x$ (règle de la chaîne).

$g'(x) = -3\sin(3x) + \sin x + x\cos x$ (règles de la chaîne et du
produit).

Règle du quotient~:
\[
h'(x) = \frac{\cos x\,(2 + \cos x) - \sin x\,(-\sin x)}{(2+\cos x)^2}
= \frac{2\cos x + \cos^2 x + \sin^2 x}{(2+\cos x)^2}
= \frac{2\cos x + 1}{(2+\cos x)^2}.
\]
\end{solution}

\begin{solution}{exo:g12:trigo:3}
\emph{(a)} Une solution particulière de $\sin x = \frac{\sqrt3}{2}$ est
$\frac\pi3$. Solutions générales~: $x = \frac\pi3 + 2k\pi$ ou
$x = \pi - \frac\pi3 + 2k\pi = \frac{2\pi}{3} + 2k\pi$. Dans
$\intoc{-\pi}{\pi}$~: $x \in \left\{\frac\pi3, \frac{2\pi}{3}\right\}$.

\emph{(b)} $\cos 2x = \frac{\sqrt2}{2} = \cos\frac\pi4$ donne
$2x = \pm\frac\pi4 + 2k\pi$, donc $x = \pm\frac\pi8 + k\pi$. Dans
$\intoc{-\pi}{\pi}$~:
$x \in \left\{-\frac{7\pi}{8}, -\frac\pi8, \frac\pi8, \frac{7\pi}{8}\right\}$.
\end{solution}

\begin{solution}{exo:g12:trigo:4}
Avec $a = \frac\pi3$, $b = \frac\pi4$~:
\[
\cos\frac{\pi}{12} = \cos a\cos b + \sin a \sin b
= \frac12\cdot\frac{\sqrt2}{2} + \frac{\sqrt3}{2}\cdot\frac{\sqrt2}{2}
= \frac{\sqrt2 + \sqrt6}{4},
\]
\[
\sin\frac{\pi}{12} = \sin a\cos b - \cos a \sin b
= \frac{\sqrt3}{2}\cdot\frac{\sqrt2}{2} - \frac12\cdot\frac{\sqrt2}{2}
= \frac{\sqrt6 - \sqrt2}{4}.
\]
\end{solution}

\begin{solution}{exo:g12:trigo:5}
Poser $s = \sin x$~: $2s^2 - s - 1 = (2s + 1)(s - 1) \geq 0$ ssi
$s \leq -\frac12$ ou $s = 1$.

$\sin x = 1$ donne $x = \frac\pi2$. $\sin x \leq -\frac12$~: sur le
cercle, l'arc sous la droite horizontale $y = -\frac12$, qui dans
$\intcc{0}{2\pi}$ est $x \in \intcc{\frac{7\pi}{6}}{\frac{11\pi}{6}}$.
D'où
\[
S = \left\{\frac{\pi}{2}\right\} \cup
\intcc{\frac{7\pi}{6}}{\frac{11\pi}{6}} .
\]
\end{solution}

\begin{solution}{exo:g12:trigo:6}
$\dfrac{\sin 3x}{x} = 3\,\dfrac{\sin 3x}{3x} \to 3 \cdot 1 = 3$
(composition avec $u = 3x \to 0$).

En utilisant $1 - \cos x = \dfrac{\sin^2 x}{1 + \cos x}$~:
\[
\frac{1 - \cos x}{x^2}
= \left(\frac{\sin x}{x}\right)^{\!2} \frac{1}{1 + \cos x}
\xrightarrow[x\to0]{} 1 \cdot \frac12 = \frac12 .
\]

En substituant $u = \frac1x \to 0^+$~:
$x \sin\frac1x = \frac{\sin u}{u} \to 1$.
\end{solution}

\begin{solution}{exo:g12:trigo:7}
$f'(x) = 1 - \sin x \geq 0$ pour tout $x$, avec égalité exactement lorsque
$\sin x = 1$, \emph{c'est-à-dire} $x = \frac\pi2 + 2k\pi$. Sur
$\intcc{0}{2\pi}$, $f$ croît de $f(0) = 1$ à $f(2\pi) = 2\pi + 1$ (le
zéro de $f'$ en $\frac{\pi}{2}$ est isolé, donc la croissance est même
stricte)~; \omterm{def:g10:functions:extrema}{minimum} $1$ en $0$, \omterm{def:g10:functions:extrema}{maximum} $2\pi + 1$ en $2\pi$. Sur $\R$,
$f' \geq 0$ partout avec seulement des zéros isolés, donc $f$ est
(strictement) \omterm{def:g12:seq:monotonic}{croissante} par le \cref{thm:g12:deriv:variations}, bien que
$f'$ s'annule en chaque $\frac\pi2 + 2k\pi$.
\end{solution}

\begin{solution}{exo:g12:trigo:8}
\emph{1.} Soit $R = \sqrt{A^2 + B^2} > 0$. Le point
$\left(\frac AR, \frac BR\right)$ est sur le cercle unité, donc il existe
$\varphi$ avec $\cos\varphi = \frac AR$ et $\sin\varphi = \frac BR$. Alors
\[
R\cos(x - \varphi) = R\bigl(\cos x\cos\varphi + \sin x \sin\varphi\bigr)
= A\cos x + B\sin x = f(x).
\]

\emph{2.} Comme $\cos$ a pour \omterm{def:g10:functions:function}{image} $\intcc{-1}{1}$, $f$ a pour \omterm{def:g10:functions:extrema}{maximum}
$R$ et pour \omterm{def:g10:functions:extrema}{minimum} $-R$. Pour $\cos x + \sin x$~: $A = B = 1$,
$R = \sqrt2$, $\varphi = \frac\pi4$, donc l'\omterm{def:g10:algebra:equation}{équation} devient
$\sqrt2\cos\left(x - \frac\pi4\right) = 1$, \emph{c'est-à-dire}
$\cos\left(x - \frac\pi4\right) = \frac{\sqrt2}{2}$, donnant
$x - \frac\pi4 = \pm\frac\pi4 + 2k\pi$~: dans $\intoc{-\pi}{\pi}$,
$x \in \left\{0, \frac\pi2\right\}$.
\end{solution}

\begin{solution}{exo:g12:trigo:9}
$(\sin)'' = (\cos)' = -\sin$ et $(\cos)'' = (-\sin)' = -\cos$~: tous deux
résolvent $y'' = -y$.

Soit $g = y - a\cos - b\sin$. Alors $g'' = -g$ (linéarité de la
dérivation), $g(0) = y(0) - a = 0$ et $g'(0) = y'(0) - b = 0$. Poser
$E = g^2 + (g')^2$. Alors
\[
E' = 2g g' + 2g' g'' = 2g g' - 2g' g = 0,
\]
donc $E$ est constante égale à $E(0) = 0$. Une somme de deux carrés qui
s'annule identiquement force $g = 0$, d'où $y = a\cos + b\sin$.
\end{solution}

\begin{solution}{pb:g12:trigo:1}
\textbf{1.} $0.99833\ldots$ et $0.99998\ldots$~: on se glisse vers
$1$.

\textbf{2.} Triangle $OAM$~: base $1$, hauteur $\sin x$, donc aire
$\frac{\sin x}{2}$. Secteur $OAM$~: la fraction $\frac{x}{2\pi}$ du
disque unité, donc aire $\frac x2$. Triangle $OAT$~: base $1$,
hauteur $\tan x$, donc aire $\frac{\tan x}{2}$.

\textbf{3.} Le triangle est contenu dans le secteur, lui-même
contenu dans le grand triangle~:
$\frac{\sin x}{2} < \frac x2 < \frac{\tan x}{2}$, c'est-à-dire
$\sin x < x < \tan x$. De $\sin x < x$ on tire
$\frac{\sin x}{x} < 1$~; et de
$x < \tan x = \frac{\sin x}{\cos x}$, en multipliant par
$\frac{\cos x}{x} > 0$~: $\cos x < \frac{\sin x}{x}$. Quand
$x \to 0^+$, $\cos x \to 1$~: encadrée, la quantité
$\frac{\sin x}{x}$ tend vers $1$~; et $\frac{\sin x}{x}$ est paire,
donc la limite bilatère vaut $1$.

\textbf{4.} $\frac{1 - \cos x}{x^2} =
\frac{1 - \cos^2 x}{x^2 (1 + \cos x)} =
\left(\frac{\sin x}{x}\right)^2 \frac{1}{1 + \cos x}
\to 1 \times \frac12$.

\textbf{5.} $\sin'(0) = \lim_{h \to 0} \frac{\sin h - 0}{h} = 1$~:
toute la dérivation du \omterm{def:g12:trigo:cossin}{sinus} et du \omterm{def:g12:trigo:cossin}{cosinus} découle de cette seule
limite. Mesurée en degrés, la limite vaudrait $\frac{\pi}{180}$, et
chaque \omterm{def:g12:deriv:derivative}{dérivée} traînerait cette constante~: le \omterm{def:g11:trigo:radian}{radian} est
exactement l'unité qui rend propre le calcul de l'oscillation.

\textbf{6.} $y' = -R\omega\sin(\omega t - \varphi)$, puis
$y'' = -R\omega^2\cos(\omega t - \varphi) = -\omega^2 y$.

\textbf{7.} $\omega = 2$. $y(0) = A = 3$~; $y'(0) = 2B = 8$, donc
$B = 4$~: $y = 3\cos 2t + 4\sin 2t$.

\textbf{8.} $R = \sqrt{9 + 16} = 5$~; période
$\frac{2\pi}{2} = \pi$~; phase
$\varphi = \arctan\frac43 \approx 0.927$~:
$y = 5\cos(2t - 0.927)$.

\textbf{9.} $5\cos(2t - \varphi) = 0$ pour la première fois lorsque
$2t - \varphi = \frac\pi2$~:
$t = \frac{\pi/2 + 0.927}{2} \approx 1.25$ s.

\textbf{10.} Avec $\sin\theta \approx \theta$~:
$\theta'' = -\frac gL \theta$, mouvement harmonique de pulsation
$\omega = \sqrt{\frac gL}$ et de période
$T = \frac{2\pi}{\omega} = 2\pi\sqrt{\frac Lg}$. Pour $L = 1$~:
$T \approx 2.006$ s. La période ne dépend pas de l'amplitude (pour
de petites oscillations) --- c'est l'isochronisme de Galilée, le
fait qui rend possibles les horloges à balancier.

\textbf{11.} $L = g\left(\frac{T}{2\pi}\right)^2 =
\frac{9.81}{\pi^2} \approx 0.994$ m~: le pendule battant la seconde
est plus court d'environ $6$ mm que le mètre du méridien.
L'Académie a choisi le méridien (la pesanteur varie d'un lieu à
l'autre, ce qui trahit le pendule)~; si $g$ avait été à un cheveu
plus grand, notre mètre battrait la seconde.

\textbf{12.} $\cos\left(x + \frac\pi2\right) =
\cos x \cos\frac\pi2 - \sin x \sin\frac\pi2 = -\sin x$~;
$\sin\left(x + \frac\pi2\right) = \cos x$. Dériver le \omterm{def:g12:trigo:cossin}{sinus} donne
donc le \omterm{def:g12:trigo:cossin}{sinus} décalé d'un quart de tour ($\cos$), puis encore
($-\sin$), puis encore, et l'on est de retour en quatre pas~:
dériver fait tourner le cercle.

\textbf{13.} En additionnant les deux développements de
$\cos(a \mp b)$~: $\cos(a-b) + \cos(a+b) = 2\cos a\cos b$. Avec
$a = b = x$~: $1 + \cos 2x = 2\cos^2 x$, c'est-à-dire
$\cos^2 x = \frac{1 + \cos 2x}{2}$.

\textbf{14.} Avec $p = \frac{p+q}{2} + \frac{p-q}{2}$ et $q$ son
symétrique, en développant puis en additionnant~:
$\sin p + \sin q = 2\sin\frac{p+q}{2}\cos\frac{p-q}{2}$. Pour
$440$ et $444$ Hz~: un son à $442$ Hz dont l'intensité est modulée
par $\cos(2\pi \cdot 2t)$ --- l'intensité culmine $4$ fois par
seconde (\omterm{def:g10:numbers:abs}{valeur absolue} de l'enveloppe)~: quatre battements par
seconde, qui s'évanouissent quand les diapasons s'accordent.

\textbf{15.} $\cos 3x = \cos 2x\cos x - \sin 2x \sin x =
(2\cos^2 x - 1)\cos x - 2\sin^2 x\cos x = 4\cos^3 x - 3\cos x$. En
$x = 20^\circ$~: $\cos 60^\circ = \frac12$, donc
$c = \cos 20^\circ$ vérifie $8c^3 - 6c - 1 = 0$ --- une \omterm{def:g10:algebra:equation}{équation} du
troisième degré sans solution par racines carrées~: construire
$20^\circ$, c'est-à-dire trisecter $60^\circ$, dépasse la règle et
le compas.

\textbf{16.} $R = \sqrt{9 + 27} = 6$~;
$\tan\varphi = \frac{3\sqrt3}{3} = \sqrt3$, donc
$\varphi = \frac\pi3$~: $I(t) = 6\cos\left(100\pi t -
\frac\pi3\right)$~: crête de $6$ ampères, déphasage de
$60^\circ$.

\textbf{17.} $\sin^2 t = \frac{1 - \cos 2t}{2}$~: période $\pi$.
$\sin 2t$ a pour période $\pi$, $\sin 3t$ a pour période
$\frac{2\pi}{3}$~: la somme se répète après le plus petit multiple
commun, soit $2\pi$.

\textbf{18.} Une oscillation de période $1$ à l'intérieur de
l'enveloppe \omterm{def:g10:functions:variations}{décroissante} $\pm\eu^{-t/10}$~: chaque balancement est
environ $10\,\%$ plus bas que le précédent. En $t = 10$,
l'enveloppe vaut $\eu^{-1} \approx 0.368$ --- encore la constante
du \cref{pb:g12:exp:1}. Les \omterm{def:g10:algebra:equation}{équations} à solutions amorties de ce
type ($y'' + a y' + b y = 0$) se résolvent au chapitre sur les
\omterm{def:g10:algebra:equation}{équations} différentielles.

\textbf{19.} Une force qui rappelle proportionnellement au
déplacement donne $y'' = -\omega^2 y$~; et
l'\cref{exo:g12:trigo:9} montre que les solutions de cette \omterm{def:g10:algebra:equation}{équation}
sont \emph{exactement} les combinaisons de $\cos$ et de $\sin$ ---
existence par exhibition, unicité par l'exercice~: la vibration n'a
pas d'autre choix que d'être sinusoïdale.

\textbf{20.} Une limite ($\frac{\sin x}{x} \to 1$, par encadrement
d'aires), deux \omterm{def:g12:deriv:derivative}{dérivées} ($\sin' = \cos$, $\cos' = -\sin$), une
\omterm{def:g10:algebra:equation}{équation} ($y'' = -\omega^2 y$), et tout un monde d'horloges, de
cordes, de courants et de marées. Astérisque~: pour de grandes
oscillations, $T = 4\sqrt{\frac Lg} \cdot
\frac{\pi/2}{M(1, \cos(\theta_0/2))}$ --- la \omterm{def:g11:stat:mean}{moyenne}
arithmético-géométrique du \cref{pb:g12:seq:1} calcule l'intégrale
elliptique en une poignée d'itérations~: la page de journal de
Gauss et le pendule de Galilée, à une formule l'une de l'autre.
\end{solution}
